\section{Leakage Controls And Shortcut Baselines}

The older cued construction is quarantined because it exposed the wrong lesson: if benchmark construction roles are visible, models may solve the interface rather than the epistemic problem. The role-uncued pilot therefore runs leakage checks before model calls and records prompt-level audits during model execution.

\begin{table}[t]
    \centering
    \small
    \caption{Leakage and prompt-hygiene gate summary for the role-uncued pilot.}
    \label{tab:leakage-gate-summary}
    \begin{tabularx}{\linewidth}{lrrX}
        \toprule
        Gate & Rows/files & Hits & Status \\
        \midrule
        Automatic leakage audit & 480 docs & 0 critical/high/medium & Passed \\
        Hidden-label scan & 480 docs & 0 hidden-label hits & Passed \\
        Semantic-ID/direct-cue scan & 480 docs & 0 semantic-ID or direct-answer hits & Passed \\
        Local surface review & 400 rows & 0 critical leaks, 0 direct cues & Passed; local leakage gate only \\
        Prompt audit & 480 prompts & 0 hidden fields, 0 audit-ID hits & Passed \\
        Stored hidden-label audit & 480 outputs & 0 prompt hits, 0 output hits & Passed \\
        \bottomrule
    \end{tabularx}
\end{table}


\Cref{tab:leakage-gate-summary} reports the main leakage controls. Automatic leakage checks found no critical, high, medium, hidden-label, semantic-ID, or direct-answer-cue hits. A local pre-model surface review covered 400 document rows and 50 unique tasks; it found no critical leaks or direct answer cues. This is not independent human-subject validation. It is a pre-model leakage gate used to decide whether model calls are justified.

\begin{table}[t]
    \centering
    \scriptsize
    \caption{Shortcut baseline summary. Values are operational epistemic escape.}
    \label{tab:shortcut-baseline-summary}
    \begin{tabularx}{\linewidth}{lrrX}
        \toprule
        Baseline family & Hidden view & Visible view & Interpretation \\
        \midrule
        ID-only & 0.000 & 0.000 & Document IDs carry no useful signal \\
        Cue phrase only & 0.000 & 0.000 & Direct cue phrases do not solve the task \\
        Metadata/source/timestamp only & 0.000 & 0.000 & Neutral metadata alone is insufficient \\
        Best non-heuristic shortcut & 0.133 & 0.133 & Below gate limit \\
        Simple heuristic & 0.350 & 0.333 & Strongest shortcut; used for model comparison \\
        \bottomrule
    \end{tabularx}
\end{table}


\Cref{tab:shortcut-baseline-summary} summarizes the no-API shortcuts. The strongest simple heuristic reaches 0.350 operational escape on the hidden view and 0.333 on the visible view. This is high enough to require caution, but low enough that model evaluation is not explained by document IDs, order, metadata alone, or a trivial cue phrase. The comparison is not uniformly favorable to every model-view pair: Claude and DeepSeek are at or below the simple heuristic on operational escape and evidence precision. The correct claim is narrower: the shortcut gates pass, and the strongest model-view pairs exceed the simple heuristic by meaningful margins.
